% s7_defects.tex -- Section 8: inert defects (main text; details in Supplementary Section S9).
% Paths in "% src:" comments are relative to the root of the repository (see README.md).
% Numbers that are not quoted verbatim from a research file are recomputed by
%   code/article/s7_defects_density_pooled.py   -> data/article/s7_defects_density_pooled.{json,csv},
%                                                      data/article/s7_defects_density_tables.txt (table rows)
%   code/article/s8_defects_placement_checks.py -> data/article/s8_defects_placement_checks.json
% Figures: code/article/s7_defects_density_figures.py (sweep), code/article/s7s8_defects_figures.py (placements).
% Convention: "\pm" is a 95 % interval (1.96 s.d./sqrt(K)) over placements unless "s.d." or "s.e." is written.
\section{Inert defects}\label{s7_defects:sec}

Throughout this section $\dd=2$, the start is the corner $\vo=(0,0)$, the target the opposite corner
$\va=(N-1,N-1)$, $q=0.8$, and $N=35$ unless stated; the clean values are $\MF=14\,100.8$ steps, $\tmode=2493$ steps and
$\tmode/\MF=0.1768$. Times are in steps. Extended tables, the validation of the codes and the proofs are in
Supplementary Section~\ref{sup_defects:sec}.
% src: data/msc_defects/structured_deterministic.csv (row clean); data/msc_defects_verify/v06_det.csv

\subsection{Blocked sites and exact computation}\label{s7_defects_density:ssec-model}

Let $\Bset\subset\Om$ be the set of blocked sites, with $\vo,\va\notin\Bset$, $M=|\Bset|$ and $p=M/N^{2}$. An attempted move
that would leave $\Om$ \emph{or land on a blocked site} is cancelled. The walker can visit only the open cluster $\Cset$,
the connected component of $\Om\setminus\Bset$ that contains $\vo$; we take $\Cset$ as the state space, write $\nn=|\Cset|$
and require $\va\in\Cset$. Open sites enclosed by blocked ones (pockets) are simply not in $\Cset$.

\begin{proposition}[Mean first-passage time on the open cluster]\label{s7_defects_density:prop-resistance}
The walk in which every move onto a blocked site or off the lattice is cancelled has a symmetric, irreducible
transition matrix on $\Cset$, and
\begin{equation}\label{s7_defects_density:eq-resistance}
  \MF_{\vo\to\va}=\frac{2\dd\,\nn}{q}\,\bigl(\Glap_{\va\va}-\Glap_{\vo\va}\bigr),
\end{equation}
where $\Glap$ is the Moore--Penrose pseudo-inverse of the graph Laplacian of $\Cset$ with unit conductances.
\end{proposition}

\begin{proof}
Two neighbouring sites of $\Cset$ exchange probability $q/(2\dd)$ per step in both directions, and every cancelled
attempt is added to the diagonal, so $I-\Pm=\frac{q}{2\dd}\Lgr_{\Cset}$ with $\Lgr_{\Cset}$ the Laplacian of $\Cset$;
Proposition~\ref{s2_model:prop-pinv} applies on $\Cset$.
\end{proof}

Blocking sites lowers the number $\nn$ of sites to be searched and, as a rule, raises the resistance factor
$\Glap_{\va\va}-\Glap_{\vo\va}$; the product need not increase. Every configuration is solved exactly: the mean by a
sparse solve, the PMF by eigendecomposition of the symmetric matrix $\Qm$ or by time stepping, and the mode by exhaustive
search together with the certificate of Proposition~\ref{s3_methods:prop-certificate}, which uses only the symmetry of
$\Qm$ and applies unchanged on $\Cset$; here it is evaluated in double precision, without the a-priori round-off bound
of Section~\ref{s3_methods:sec} (Supplementary Section~\ref{sup_defects:ssec-validation}). Two implementations, written separately and sharing no code, agree with each
other and with a direct simulation of the walker rule (Table~\ref{s7_defects_density:tab-validation}); they produced two
independent samples of placements, \emph{sample~A} and \emph{sample~B}.

\subsection{Uniform and corner-thinned placement}\label{s7_defects_density:ssec-sweep}

We solved $400$ placements at each of $13$ densities, $p=0.02$ to $0.35$, under two placement laws, in each of the two
samples ($20\,800$ configurations in all): \emph{uniform} (every site except $\vo$ and $\va$ equally likely) and
\emph{corner-thinned}, in which $M=\mathrm{round}(pN^{2})$ sites are sampled without replacement with weight $1$ except in
the $4\times4$ block of sites at each of the two corners, where the weight is $\exp[-(4-\rho)/2]$ with $\rho$ the
Euclidean distance to that corner ($\eu^{-3/2}$ at the two neighbours of the corner, rising to $\eu^{0.12}\approx1.13$ at the far
corner of the block, where $\rho=3\sqrt2$; the corner itself is never blocked). It is a simple placement law correlated with the positions of start and target: it
thins the defects next to them and leaves the bulk uniform. Placements are redrawn until $\va\in\Cset$. The two samples agree: over the $26$ cells the
largest standardised difference between them is $1.9$ for the mean, $2.1$ for the mode and $1.9$ for the ratio.
% src: data/msc_defects_verify/v03_summary.csv (z_mfpt, z_mode, z_ratio); data/article/s7_defects_density_pooled.csv (z_ratio_A_vs_B); code/msc_defects/ensemble.py (generate_defects_smart: corner-thinned weights)
Figure~\ref{s7_defects_density:fig-sweep} shows the pooled results; Table~\ref{s7_defects_density:tab-sweep} gives them in
full. Two estimators of the ratio must be kept apart, the mean of the per-placement ratios, $\langle\tmode/\MF\rangle$, and
the ratio of ensemble means, $\langle\tmode\rangle/\langle\MF\rangle$; between placements $\MF$ fluctuates more than $\tmode$
(relative standard deviations $26\%$ and $10\%$ at $p=0.2$ under uniform placement), and the first estimator lies above
the second in every cell.
% src: data/article/s7_defects_density_pooled.json (statements.uniform_relsd_mfpt_and_mode_by_p, mean_of_ratios_above_ratio_of_means_in_all_cells)

\emph{Uniform placement.} Up to $p=0.14$ the ratio is unchanged to within a few per cent, with a sign that depends on the
estimator: $\langle\tmode\rangle/\langle\MF\rangle$ lies between $0.8\%$ below and $0.3\%$ above the clean value, whereas
$\langle\tmode/\MF\rangle$ exceeds it by at most $2.4\%$ for $p\le0.12$ and by $3.2\%$ at $p=0.14$; between $27\%$ and $37\%$ of the
individual placements with $p\le0.2$ have a ratio \emph{below} the clean one. From $p=0.16$ both estimators rise, to
$0.1926\pm0.0022$ ($+8.9\%$) and $0.1851$ ($+4.7\%$) at $p=0.2$. At $p=0.2$ the mean and the mode are multiplied by
$1.825\pm0.033$ and $1.911\pm0.013$, while the coefficient of variation hardly moves ($0.918\to0.910$): the distribution is
delayed, not reshaped.
% src: data/article/s7_defects_density_pooled.json (statements.uniform_rom_excess_pct_by_p, uniform_ratio_excess_pct_by_p, uniform_frac_below_clean_range_p_le_0.20); data/article/s7_defects_density_pooled.csv (mfpt_fac_P, mode_fac_P, cv_P at p = 0.20)

\emph{Corner-thinned placement.} Here $\langle\tmode/\MF\rangle$ increases monotonically with $p$: $0.1794\pm0.0003$,
$0.1904\pm0.0010$, $0.2066\pm0.0020$ and $0.2550\pm0.0050$ at $p=0.02$, $0.10$, $0.20$ and $0.35$. A least-squares line through
the per-placement ratios for $p\le0.2$ has slope $+0.159\pm0.007$ per unit $p$, against $+0.081\pm0.009$ under uniform
placement. At $p=0.2$ the mean and the mode are multiplied by $1.606\pm0.023$ and $1.829\pm0.011$. Section~\ref{s8_defects_placement:ssec-capture} explains the difference: thinning the defects next to the target, where the mean is most
sensitive, slows the mean less than the mode.
% src: data/article/s7_defects_density_pooled.json (sweep; statements.corner_thinned_pooled_ratio_monotone_in_p; ols_trend_ratio_vs_p.corner_thinned_pmax0.20_P, uniform_pmax0.20_P); data/article/s7_defects_density_pooled.csv (mfpt_fac_P, mode_fac_P at p = 0.20)

\begin{figure}[tbp]
\centering
\includegraphics[width=\textwidth]{s7_defects_density_sweep.pdf}
% src: code/article/s7_defects_density_figures.py; data data/article/s7_defects_density_pooled.csv, data/msc_defects/homogenization.json
\caption{Density sweep, $N=35$, $q=0.8$, corner to opposite corner, discrete time. (a)~Mean and (b)~mode divided by
their clean values, and (c)~$\tmode/\MF$, against the blocked fraction $p$, for uniform placement (circles) and
corner-thinned placement (squares). Filled symbols: sample~A; open symbols: sample~B ($400$ placements each; bars are
$95\%$ intervals over placements; the samples are displaced horizontally for clarity). In (a) and (b) the solid curve
is the homogenised factor $\Thetap_{192}(p)$ measured on tori (Table~\ref{s7_defects_density:tab-homog}) and the dashed
curve the dilute law \eqref{s7_defects_density:eq-dilute} in the form $1/[1-(\pi-1)p]$. In (c) the symbols are means of
per-placement ratios and the lines the pooled ratio of ensemble means; the horizontal line is the clean value $0.1768$.}
\label{s7_defects_density:fig-sweep}
\end{figure}

\subsection{Effective medium and time dilation}\label{s7_defects_density:ssec-emt}

The walk with cancelled moves on a site-diluted lattice is the ``blind ant'' of the percolation literature
\citep{HavlinBenAvraham1987}. Its stationary distribution is uniform on the accessible cluster, and the Einstein relation
ties the long-distance diffusivity $\Deff$ to the conductivity $\cond$ of the network of unit resistors on the open
bonds: $\cond/\condz=\Pinf\,\Deff/\Dzero$, where $\Pinf$ is the fraction of sites in the accessible cluster
\citep{HavlinBenAvraham1987,StaufferAharony1994}. We define the time-dilation factor
\begin{equation}\label{s7_defects_density:eq-theta}
  \Thetap(p)=\frac{\Dzero}{\Deff(p)}=\Pinf(p)\,\frac{\condz}{\cond(p)} .
\end{equation}
The hypothesis to be tested is that, on scales large compared with the percolation correlation length, uniformly random
defects multiply every time scale by $\Thetap(p)$ and leave the shape of the first-passage law, hence $\ratio$, unchanged.
For the mean this is the coarse-grained form of \eqref{s7_defects_density:eq-resistance}. $\Thetap$ diverges at the blocked
fraction $p_{\mathrm c}=1-0.5927=0.4073$ \citep{NewmanZiff2000}.
% src: Newman and Ziff (2000): site-percolation threshold of the square lattice 0.59274621(13)

At small $p$ the conductivity follows from the response of a single blocked site. Let $\pk$ be the potential kernel of
the planar simple random walk \citep{Spitzer1976}, with $\pk(\mathbf 0)=0$, $\pk(\pm\ve_i)=1$ and $\pk(\pm2\ve_i)=4-8/\pi$.
% src: Spitzer1976 (potential kernel; values recomputed by quadrature in data/article/s7_defects_density_dipole_check.json: a10 = 1.0000000000000002, a20 = 1.4535209105296745 = 4 - 8/pi)

\begin{proposition}[Dipole response of one blocked site]\label{s7_defects_density:prop-dilute}
On the infinite square lattice with unit conductances, remove the site $\mathbf 0$ together with its four bonds and
impose the potential $-E x_1$ at infinity. Then:
(a) the potential is unique, $U(\vx)=-Ex_1-\frac{\pi E}{8}\,[\pk(\vx+\ve_1)-\pk(\vx-\ve_1)]$, and at the two neighbours
$\pm\ve_1$ of the removed site it is $\mp\pi E/2$ (instead of $\mp E$);
(b) the removed site lowers the dissipated power by $\pi E^{2}$, in the sense that $\sum_{b}U^{0}_{b}U_{b}=\pi E^{2}$ over
the four removed bonds, and in every finite network of unit resistors with prescribed boundary potentials removing a
set of bonds lowers the dissipated power by exactly $\sum_{b\ \mathrm{removed}}U^{0}_{b}U_{b}$.
\end{proposition}

\begin{heuristic}[Dilute law]\label{s7_defects_density:res-dilute}
For uniformly placed blocked sites on the square lattice, as $p\to0$,
\begin{equation}\label{s7_defects_density:eq-dilute}
  \frac{\cond}{\condz}=1-\pi p+\Order(p^{2}),\qquad
  \frac{\Deff}{\Dzero}=1-(\pi-1)\,p+\Order(p^{2}),\qquad
  \Thetap(p)=1+(\pi-1)\,p+\Order(p^{2}).
\end{equation}
\end{heuristic}

The derivation treats each blocked site as isolated in the unperturbed field (Proposition~\ref{s7_defects_density:prop-dilute}
and $\Pinf=1-p+\Order(p^{4})$); the step from one defect to a density of defects is not proved here, hence the label. The
law is not new: \citet{WatsonLeath1974} gave $\cond/\condz=1-\pi p+\tfrac12\pi p^{2}$ from an effective-medium theory, and
\citet{Nieuwenhuizen1986} calculated the diffusion coefficient and conductivity of exactly this walk to order $p^{2}$, with
the coefficients printed in \citet{Ernst1987}:
\begin{equation}\label{s7_defects_density:eq-second-order}
  \frac{\cond}{\condz}=1-\pi p+1.2858\,p^{2}+\cdots,\qquad
  \frac{\Deff}{\Dzero}=1-(\pi-1)\,p-0.8558\,p^{2}+\cdots .
\end{equation}
Our own evaluation of the pair term on tori gives $1.28588$ and $-0.85571$, within $10^{-4}$ of the printed values, and
the conductivities measured on $192\times192$ tori agree with \eqref{s7_defects_density:eq-second-order} to $6.2\times10^{-4}$ for
$p\le0.18$; the estimates $\Thetap_{192}$ of the time-dilation factor are $1.045$, $1.289$, $1.877$ and $3.739$ at $p=0.02$, $0.1$,
$0.2$ and $0.3$ (Supplementary Table~\ref{s7_defects_density:tab-homog}).
% src: data/article/s7_defects_density_second_order.json (a2_extrapolated_linear_in_Nt^-2_last3, D2_coefficient_from_extrapolated_a2, extrapolated_minus_published); data/article/derived_checks.json (E_dilute_laws: max_abs_sigma_minus_exact_second_order_p_le_0.18); data/msc_defects/homogenization.json (finite_p); coefficients 1.2858 and -0.8558: Ernst, Nieuwenhuizen and van Velthoven (1987), Eqs (4.2), (4.3)

\emph{Test at $N=35$.} Under uniform placement and $p\le0.2$ the pooled factor of the mean is within $2.8\%$ of
$\Thetap_{192}(p)$ and that of the mode within $1.82\%$, with no fitted parameter (Fig.~\ref{s7_defects_density:fig-sweep}(a,b));
these are ensemble-mean estimates, whose $95\%$ intervals reach $\pm1.8\%$ for the mean and $\pm0.72\%$ for the mode
(relative to $\Thetap_{192}$), so that the deviations of the mean are of the size of their sampling uncertainty.
For $p\ge0.25$ the mean falls below $\Thetap$ ($3.08$ against $3.74$ at $p=0.3$) while the mode still follows it ($3.66$);
those densities describe a selected ensemble on a lattice not much larger than the correlation length. The simplest
mean-field description of defects, a reduced activity $q_{\mathrm{eff}}(p)$, cannot capture the changes of $\ratio$ found
below: a change of activity only rescales time, and by \eqref{s2_model:eq-qscaling} and
Proposition~\ref{s2_model:prop-structure}(d) it leaves $\ratio$ unchanged up to a few steps
(Section~\ref{s6_mechanism:sec}).
% src: data/article/s7_defects_density_pooled.json (statements.uniform_max_abs_mfpt_fac_over_theta_pct_p_le_0.20_pooled, uniform_max_abs_mode_fac_over_theta_pct_p_le_0.20_pooled); data/article/s7_defects_density_pooled.csv (uniform, p <= 0.20: mfpt_fac_ci_P / theta_L192 up to 1.77 %, mode_fac_ci_P / theta_L192 up to 0.719 %)

\emph{Larger lattices.} For uniform placement at $p=0.1$ and $0.2$ we computed modes up to $N=200$, certified by
Proposition~\ref{s3_methods:prop-certificate} in double precision, and means up to $N=280$, with the most precise estimates $\Thetap_{200}(0.1)=1.2877$ and $\Thetap_{200}(0.2)=1.8783$. The mode factor equals
$\Thetap$ to within $1.11\%$ at $p=0.1$ and $1.6\%$ at $p=0.2$ for every $N$ from $15$ to $200$; the factor of the mean divided by
$\Thetap$, averaged over $N=35$--$280$, is $1.004\pm0.005$ and $0.987\pm0.007$. The shift of $\langle\tmode/\MF\rangle$ relative
to the clean lattice of the same size is between $+0.6\%$ and $+1.2\%$ for $N\ge50$ at $p=0.1$; at $p=0.2$ it falls from
$+9.4\pm1.8\%$ at $N=35$ (in the $400$ placements of this size study; the $800$ pooled placements of Section~\ref{s7_defects_density:ssec-sweep} give $+8.9\%$) to $+4.5\pm1.4\%$ at $N=100$ and $+3.9\pm2.7\%$ at $N=200$: it halves but is not shown to vanish
(Supplementary Table~\ref{s7_defects_density:tab-size}). A slow approach is to be expected, because in
\eqref{s7_defects_density:eq-resistance} the constant that accompanies $\ln N$ in $\Glap_{\va\va}$ is set by the immediate
surroundings of the target.
% src: data/article/s7_defects_density_pooled.json (size_dependence.summary); data/msc_defects_verify/v07_summary.csv, v07b_summary.csv, v04b_theta.json

\begin{conjecture}[Homogenisation]\label{s7_defects_density:conj-homog}
Fix $p<p_{\mathrm c}$ and place $\mathrm{round}(pN^{2})$ blocked sites uniformly, conditioned on $\va\in\Cset$. As
$N\to\infty$, $\langle\MF\rangle/(\Thetap(p)\,\MF_{\mathrm{clean}})\to1$ and
$\langle\tmode\rangle/(\Thetap(p)\,\tmode_{\mathrm{clean}})\to1$, so that $\langle\tmode\rangle/\langle\MF\rangle$ divided by the
clean ratio at the same $N$ tends to one.
\end{conjecture}

\subsection{Placement decides the direction}\label{s8_defects_placement:ssec-capture}

By \eqref{s7_defects_density:eq-resistance}, $\MF=(4\nn/q)(\Glap_{\va\va}-\Glap_{\vo\va})$. The rows of $\Glap$ sum to zero, so
$(4\nn/q)\,\Glap_{\va\va}$ is the mean first-passage time from a uniformly distributed start, and the actual start enters
only through the cross term $\Glap_{\vo\va}$ ($-0.221$ against $\Glap_{\va\va}=2.082$ on the clean lattice). We call
$\Glap_{\va\va}$ the \emph{capture resistance} of the target. It controls the mean: over the 400 uniform placements at
$p=0.1$ and at $p=0.2$ the rank correlation between $\MF$ and $\Glap_{\va\va}$ is $0.9995$ and $0.9986$ (sample~A; $0.9995$ and
$0.9982$ in sample~B), whereas that between $\MF$ and the number of blocked sites near the start is below $0.08$ in
absolute value; and at $p=0.2$ the uniform and the corner-thinned placements leave the same accessible volume
($\nn/N^{2}=0.797$) but raise $\Glap_{\va\va}-\Glap_{\vo\va}$ by the factors $2.310$ and $2.022$.
% src: data/msc_defects/structured_deterministic.csv (row clean: G_tt, G_st); data/msc_defects/summary_correlations.csv (G_tt, blk_s_r3, blk_s_r6); data/msc_defects_verify/v09_mechanism.json (G_aa, blk_s6); data/msc_defects/summary_sweep_N35.csv (n_factor, Gdiff_factor)
The mode responds to a different part of the lattice. Blocking a single site shortens the mean for $67.6\%$ of the sites
(the lost volume), and the positive responses are concentrated around the target, where $64\%$ of $\sum_{\vx}\chiT(\vx)$ comes
from sites within distance 3; the response of the mode is spread over the whole region between start and target
($19\%$ within distance 3 of the target), and only $5.4\%$ of the sites lower it
(Fig.~\ref{s8_defects_placement:fig-maps}).
% src: data/msc_defects/summary_first_order.json; data/msc_defects_verify/v05_sens_N35.json; data/article/s8_defects_placement_checks.json (susceptibility_N35)

\begin{table}[tbp]
\centering
\caption{Localised random placements of $M$ blocked sites ($N=35$, $q=0.8$, 200 placements per row and sample,
mean $\pm$ 95\% confidence interval). Factors are relative to the clean lattice, for which $\tmode/\MF=0.1768$.}
\label{s8_defects_placement:tab-local}
\small
\begin{tabular}{lrcccc}
\toprule
 & & \multicolumn{3}{c}{sample A} & sample B\\
\cmidrule(lr){3-5}\cmidrule(lr){6-6}
region & $M$ & $\MF$ factor & $\tmode$ factor & $\tmode/\MF$ & $\tmode/\MF$\\
\midrule
anywhere    & 40 & $1.086\pm0.017$   & $1.080\pm0.005$ & $0.1771\pm0.0016$ & $0.1780\pm0.0016$\\
centre      & 40 & $1.0086\pm0.0008$ & $1.150\pm0.003$ & $0.2015\pm0.0004$ & $0.2014\pm0.0004$\\
near start  & 10 & $0.9933\pm0.0001$ & $1.012\pm0.001$ & $0.1801\pm0.0001$ & $0.1802\pm0.0002$\\
near start  & 40 & $0.999\pm0.003$   & $1.320\pm0.027$ & $0.2329\pm0.0040$ & $0.2368\pm0.0048$\\
near target & 10 & $1.239\pm0.029$   & $1.056\pm0.006$ & $0.1530\pm0.0020$ & $0.1530\pm0.0019$\\
near target & 20 & $1.659\pm0.057$   & $1.134\pm0.009$ & $0.1249\pm0.0024$ & $0.1276\pm0.0025$\\
near target & 30 & $2.281\pm0.083$   & $1.230\pm0.011$ & $0.0994\pm0.0023$ & $0.0991\pm0.0024$\\
near target & 40 & $3.367\pm0.127$   & $1.357\pm0.015$ & $0.0752\pm0.0022$ & $0.0741\pm0.0021$\\
\bottomrule
\end{tabular}
% src: data/msc_defects/summary_structured_local.csv (sample A); data/msc_defects_verify/v06_local.csv via data/article/s8_defects_placement_checks.json (local_sample_B)
\end{table}

To separate placement from density we block $M$ sites drawn uniformly inside one region: within Euclidean distance 12 of
the target (122 candidate sites), within distance 12 of the start (122 sites), within $6.2$ of the centre (121 sites), or
anywhere (Table~\ref{s8_defects_placement:tab-local}, Fig.~\ref{s8_defects_placement:fig-structured}). Forty blocked sites,
$3.3\%$ of the lattice, multiply the mean by $3.37\pm0.13$ and the mode by $1.36$ near the target, so that the ratio falls to
$0.0752\pm0.0022$, less than half its clean value; near the start they leave the mean unchanged ($\times0.999\pm0.003$),
delay the mode by a third and raise the ratio to $0.2329\pm0.0040$. Deterministic obstacles behave in the same way: serpentine
corridors push the law towards the one-dimensional shape ($\tmode/\MF=0.316$--$0.339$), walls between start and target
raise the ratio to between $0.244$ and $0.328$, whereas a wall next to the target lowers it to $0.135$; enclosing the
target makes the law more nearly exponential (ratio $0.0925$ for ten blocked sites), and a solid $12\times12$ block in the centre
\emph{lowers} the mean by $4.6\%$ and delays the mode by $23\%$. Hence $\tmode/\MF$ is not a function of the defect density:
it decreases when the defects surround the target, increases when they lie around the start or between start and target,
and, when they are spread uniformly, stays within a few per cent of its clean value at low density ($p\le0.14$;
Section~\ref{s7_defects_density:ssec-sweep}) but rises by
$4.7\%$ to $8.9\%$, depending on the estimator, at $p=0.2$.
% src: data/article/s7_defects_density_pooled.json (statements.uniform_ratio_excess_pct_by_p, uniform_rom_excess_pct_by_p: p = 0.14 and 0.20)
% src: data/msc_defects/summary_structured_local.csv (mfpt_factor, mode_factor, ratio); data/msc_defects/structured_deterministic.csv via data/article/s8_defects_placement_checks.json (deterministic, deterministic.block12_centre)

\begin{figure}[tbp]
\centering
\includegraphics[width=\textwidth]{s8_defects_placement_structured.pdf}
\caption{Placement decides the direction ($N=35$, $q=0.8$, corner to opposite corner, discrete time).
(a)~Mode factor against mean factor for the 28 deterministic geometries of
Table~\ref{appC_tables:tab-deterministic}, logarithmic axes; on the solid line $\tmode/\MF$ keeps its clean value,
on the dashed lines it is doubled or halved. (b)~Mean (solid lines) and mode (dashed lines) relative to the clean
lattice for $M$ blocked sites drawn inside a region. (c)~The resulting ratio; the grey line is the clean value.
In (b) and (c) the points are means over the 200 placements of sample~A and the bars 95\% confidence intervals.}
% src: data/msc_defects/structured_deterministic.csv, summary_structured_local.csv (figure: code/article/s7s8_defects_figures.py, fig_structured)
\label{s8_defects_placement:fig-structured}
\end{figure}

\emph{Two time scales.} Beyond the transit time the PMF is dominated by its two slowest terms, with time constants
$\Tcap=-1/\ln(1-q\nu_0)$ (the capture time) and $\tau_1=-1/\ln(1-q\nu_1)$; with $1/\tau=1/\tau_1-1/\Tcap$ and
$c=-b_1/b_0$, the maximiser of the two-term sum is $\tmode\approx\tau\ln[c(1+\Tcap/\tau)]$. A pure dilation multiplies $\Tcap$
and $\tau$ by the same factor and leaves $\tmode/\MF$ unchanged, whereas a perturbation that raises only the capture time
moves the mode sub-linearly: on the clean lattice the exponent $\partial\ln\tmode/\partial\ln\Tcap$ is $0.21$ with $\tau$
and $c$ held fixed, and $0.17$ with $\tau_1$ and $c$ held fixed ($\tau$ then follows $\Tcap$). The exact
elasticities $\ln(\tmode\ \text{factor})/\ln(\MF\ \text{factor})$ of obstacles that enclose the target are somewhat larger,
$0.234$ to $0.310$, and the regression of $\ln\tmode$ on $\ln\MF$ across the uniform placements at $p=0.2$ has slope $0.30$. This
heuristic accounts for the sign and the approximate size of the effect and has no controlled error; it fails for a wall
near the target, which also lengthens the transit (Supplementary Section~\ref{sup_defects:ssec-twotime}).
% src: data/article/s8_defects_placement_checks.json (two_time.elasticity_formula, two_time.elasticity_formula_tau1_c_fixed, two_time.cross_placement_slope_p02); data/msc_defects_verify/v09b_twotime.csv (elasticity_exact); data/msc_defects_verify/v09_mechanism.json (p0.20.loglog_slope_mode_on_mfpt)

\emph{A one-parameter family of shapes.} Over all $13\,631$ exactly solved configurations of sample~A the coefficient of
variation is nearly a function of $\ratio$ alone (Pearson correlation $-0.981$; a quadratic fits with a root-mean-square
residual of $0.0045$); sample~B gives $-0.983$. We record this as an empirical regularity, without a theory: size, density
and placement appear to act mainly by moving the law along one line between a transit-limited shape and a
capture-limited, nearly exponential one.
% src: data/article/s8_defects_placement_checks.json (shape_family.sample_A); data/msc_defects_verify/v09_mechanism.json (family_main_implementation, family_recheck)

\emph{Other defect types.} On a chain, a barrier crossed with reduced probability costs more the nearer it is to the target
(Proposition~\ref{s8_defects_placement:prop-1d-barriers}; for a single barrier this is Eq.~(12) of
\citet{SarvaharmanGiuggioli2023}, who also showed that a barrier behind the start changes the mode and the tail but not
the mean, their `disorder indifference'), and uniformly random barriers are again a dilation of time. On the
square lattice, permeable obstacles at $p=0.1$ leave the ratio within $3.2\%$ of its clean value, but the limit of vanishing
permeability is singular, because a site that can be entered at all carries its full share of the stationary distribution.
Partial traps, which remove the walker with probability $\trap$ per step, remove the long trajectories and act in the
opposite sense to obstacles: the ratio of mode to conditional mean rises from $0.178$ to $0.638$ as $\trap$ grows from $10^{-5}$
to $10^{-2}$ at $p=0.1$ (Supplementary Section~\ref{sup_defects:ssec-other}).
% src: data/msc_defects/summary_beyond.json (2d_permeable, 2d_traps); data/msc_defects_verify/v08_beyond.json (perm, traps); data/article/s8_defects_placement_checks.json (permeable)

Within this model, and at the sizes computed, the mean and the mode respond to different parts of the lattice: the mean
to the capture resistance of the target, which is set by its immediate surroundings, and the mode to the transit between
start and target. We make no claim beyond the lattices, sizes and defect models examined here.
